\subsection{A Graph Method}
\label{sec:graph_method}

In the earlier subsection it was shown that even though one starts from an extended system, for the evaluation of the higher order asymptotic terms of the center-manifold only  the standard linear system is required and the system extension only affects the tangent (linear) space. This aspect however is not unique to the normal form and other parameterizations can be constructed in which the asymptotic center-manifold of a system can be evaluated using just the standard linear operator $\Lu$. Here we construct a particular graph style parameterization.

Starting again with the asymptotic representations of $\Param(\xbf)$ and $\mathcal{Y}(\xbf)$ as in equations~\eqref{eqn:reduced_normal_form} and \eqref{eqn:asymptotic_normal_form} respectively, one may constrain all higher order terms $\yhatbf_{i_{1},\ldots,i_{k}}$ to lie in the complement of the center subspace (spanned by $\Vhat$), resulting in,
\begin{align}
	\label{eqn:graph_condition}
	\whatbf^{\dagger}_{j}\yhatbf_{i_{1},\ldots,i_{k}} = \veczero, && \forall j\in1,\ldots,m,
\end{align}
for all $(i_{1},\ldots,i_{k})\in\IndexSet{k}{1,m},\ k\ge2$. Note that this condition applies regardless of resonance. 

At $\mathcal{O}(|\xbf|)$ one again obtains equation~\eqref{eqn:asymptotic_O1} and as in the case of the normal form, the unknowns $\yhatbf_{i_{1}},\AsParam_{i_{1}}$ \etc, can be associated with the generalized eigenvalue problem resulting in equation~\eqref{eqn:GEV_O1}. At first order the expressions are identical to those obtained for the normal form, which is unsurprising since the $\mathcal{O}(|\xbf|)$ terms only represent the center subspace.

One may therefore represent the general expression for the evaluation of higher order terms as,
\begin{align}
	\label{eqn:graph_Ok_compact}
	\begin{aligned}
		\left[\omega_{i_{1},\ldots,i_{k}}\Ihat - \ExtL\right]\yhatbf_{i_{1},\ldots,i_{k}} + \sum^{m}_{j=1}\AsParamEl^{j}_{i_{1},\ldots,i_{k}}\vhatbf_{j} \\
		= \fhatbf_{i_{1},\ldots,i_{k}} - \hhatbf^{O}_{i_{1},\ldots,i_{k}} 
		- \hhatbf^{L}_{i_{1},\ldots,i_{k}}.
	\end{aligned} &&
\end{align}
Since the graph parameterization is defined by equation~\eqref{eqn:graph_condition}, one obtains the definition of the asymptotic terms of $\Param$ as,
 \begin{align}
 	\label{eqn:graph_Ok_param_compact}
 		 \AsParamEl^{j}_{i_{1},\ldots,i_{k}} 
 		= \whatbf^{\dagger}_{j}(\fhatbf_{i_{1},\ldots,i_{k}} - \hhatbf^{O}_{i_{1},\ldots,i_{k}} 
 		- \hhatbf^{L}_{i_{1},\ldots,i_{k}}),
 \end{align}
$\forall j\in1,\ldots,m$. The extended variables vanish for $\fhatbf_{i_{1},\ldots,i_{k}}$ by construction and following the logical arguments presented in the previous section, one may again show that the extended variables also vanish for $\hhatbf^{O}_{i_{1},\ldots,i_{k}}, \hhatbf^{L}_{i_{1},\ldots,i_{k}}$ and $\yhatbf_{i_{1},\ldots,i_{k}}$. This then implies that, 
\begin{align}
	\label{eqn:graph_Ok_param_compact2}
	\AsParamEl^{j}_{i_{1},\ldots,i_{k}}= 0, && \forall j\in n+1,\ldots,m.
\end{align}


%A particular graph-style parameterization has been used in \cite{negi24} for the construction of center-manifolds of the Navier--Stokes while treating variations of the Reynolds number as parameter perturbations. 


%Since the parameterization can be chosen arbitrarily, one can always choose a representation in which $\mathcal{Y}(\xbf)$ has non trivial components in the extended variables even for the higher order asymptotic terms $\yhatbf_{i_{1},\ldots,i_{k}}$.

%One may find the exposition on the implementation of the parameterization method in previous works of \citet{haro16} and \citet{jain22}. To be clear, the application of the parameterization method to an extended system like equation~\eqref{eqn:evolution_extended} follows the same steps as outlined in those earlier works. However, many of the essential steps are repeated here to bring out the structure of system of equations at each asymptotic order in the context of system extension.



















